\section{Related Work}
\label{sec:related}

\para{Comparing tokenizers.}
BPE started as a compression algorithm~\citep{gage1994new} and was adapted to subword segmentation for neural machine translation~\citep{sennrich2016neural}.
GPT-2 later made it byte-level~\citep{radford2019language}.
Several studies compare tokenizer families.
\citet{bostrom2020byte} find that unigram language-model tokenization~\citep{kudo2018subword} matches or beats BPE for pre-training.
\citet{ali2024tokenizer} show that tokenizer choice affects the downstream performance and training cost of multilingual LLMs.
\citet{schmidt2024tokenization} argue that tokenization is more than compression, because tokenizers with similar compression can give different downstream results.
These studies vary the algorithm, the vocabulary, and often the training data together.
Unlike them, we fix the algorithm, the corpus, and the vocabulary size, and we vary only merge order.
Our same-vocabulary controls go one step further and fix the vocabulary itself.

\para{Compression and intrinsic measures.}
\citet{galle2019investigating} link the success of BPE to shorter sequences.
\citet{goldman2024unpacking} show that compression correlates with downstream quality, especially for generation.
\citet{zouhar2023tokenization} propose R\'enyi efficiency as a better predictor of translation quality than compression.
We test both measures on tokenizers that differ only in merge order.
Bytes-per-token rank-predicts \bpb perfectly in our setting, while R\'enyi efficiency is a much weaker predictor.
We also find that the link between compression and quality holds for language modeling but not for fine-tuned classification.

\para{Stochastic and alternative segmentation.}
BPE-dropout~\citep{provilkov2020bpe} and subword regularization~\citep{kudo2018subword} randomly perturb segmentation during training to improve robustness.
These methods change segmentation at training time while keeping the standard segmentation at test time.
Our \randorder and \revorder tokenizers instead fix a different, deterministic merge priority for both training and testing.
At the extreme, our \revlearn tokenizer becomes almost character-level, which connects to byte-level models such as ByT5~\citep{xue2022byt5}.
Unlike ByT5, we keep the model and context length fixed, so a shorter-unit tokenizer sees less text per window.

\para{Domain-specific vocabularies.}
Domain models such as SciBERT~\citep{beltagy2019scibert} and PubMedBERT~\citep{gu2021domain} build vocabularies from in-domain text, and \citet{gururangan2020don} show that continued domain pre-training helps.
These works change the vocabulary and the pre-training data together.
Our \domain tokenizer changes only merge order: it up-weights biomedical counts in the same mixed corpus.
This isolates how much of a domain gain comes from the tokenizer alone.
